\begin{textsample}{NEG Dim 5 – global\_north\_2024\_07 – Score -1.00 – global\_north\_2024\_07/110189300.txt}  \label{ex:f5_neg_154}
US-built pier will be put back in Gaza for several days to move aid, then permanently removed WASHINGTON ( AP ) -- The pier built by the U.S. military to bring humanitarian aid to Gaza will be reinstalled Wednesday to be used for several days, but then the plan is to pull it out permanently, several U.S. officials said.
It would deal the final blow to a project long plagued by bad weather, security uncertainties and difficulties getting food into the hands of starving Palestinians.
The officials said the goal is to clear whatever aid has piled up in Cyprus and on the floating dock offshore and get it to the secure area on the \textbf{beach} in Gaza.
Once that has been done, the Army will dismantle the pier and depart.
The officials spoke on condition of anonymity because final details are still being worked out.
Officials had hoped the pier would provide a critical flow of aid to starving residents in Gaza as the nine-month-long war drags on.
But while more than 19.4 million pounds ( 8.6 million kilograms ) of food has @ @ @ @ @ @ @ @ @ @ been hampered by persistent heavy seas and stalled deliveries due to ongoing security threats as Israeli troops continue their offensive against Hamas in Gaza.
The decision comes as Israeli troops make another push deeper into Gaza City, which Hamas says could threaten long-running negotiations over a cease-fire and hostage release, after the two sides had appeared to have narrowed the gaps in recent days.
U.S. troops removed the pier on June 28 because of bad weather and moved it to the port of Ashdod in Israel.
But distribution of the aid had already stopped due to the security concerns.
The United Nations suspended deliveries from the pier on June 9, a day after the Israeli military used the area around it for airlifts after a hostage rescue that killed more than 270 Palestinians.
U.S. and Israeli officials said no part of the pier itself was used in the raid, but U.N. officials said any perception in Gaza that the project was used may endanger their aid work.
As a result, aid brought through the @ @ @ @ @ @ @ @ @ @ for days while talks continued between the U.N. and Israel.
More recently, the World Food Program hired a contractor to move the aid from the \textbf{beach} to prevent the food and other supplies from spoiling.
The Pentagon said all along that the pier was only a temporary project, designed to prod Israel into opening and allowing aid to flow better through land routes -- which are far more productive than the U.S.-led sea route.
And the weather now is projected only to get worse.
The pier was damaged by high winds and heavy seas on May 25, just a bit more than a week after it began operating, and was removed for repairs.
It was reconnected on June 7, but removed again due to bad weather on June 14.
It was put back days later, but heavy seas again forced its removal on June 28.

% matched lemmas: beach
\end{textsample}
